Then the ordinary system of equation with 
$\partial /\partial  \gamma=d/d \gamma=()'$ 

\begin{equation} 
F'''
+
F
F''
=0
\end{equation} 

with the boundary condition:
$F(0):=F_0=0$ (v=0 at the dividing streamline), $F'(\infty)=F_\infty=\dfrac{u_{\infty}}{u_{-\infty}}=1$ and  $F'(-\infty):=F'_{-\infty}=\dfrac{u_{-\infty}}{u_{\infty}}=0.5$.

Remember that F is self-similar equation under the invariance to the transformation $F(\gamma)=a\tilde{F}(a\gamma)$. 

To corroborate it, define $\upsilon=a\gamma$ which imply
$ 
F(\gamma)=a\tilde{F}(\upsilon)$,
then using the chain ruler  

\begin{equation*}
F'(\gamma)= 
\frac{d F}{ d\gamma }
=
\frac{d F}{ d\upsilon }
\frac{d\upsilon }{d\gamma}
=
a
\frac{d (a\tilde{F})}{ d\upsilon }
=
a^2
\frac{d \tilde{F}}{ d\upsilon }
=
a^2
\tilde{F}'
\end{equation*}

From de second derivative:

\begin{equation*}
F''(\gamma)= 
\frac{d }{ d\gamma }
\left[
\frac{d F}{ d\gamma }
\right]
=
\frac{d }{ d\upsilon }
\left[
\frac{d F}{ d\gamma }
\right]
\frac{d\upsilon }{d\gamma}
=
a
\frac{d }{ d\upsilon }
\left[
\frac{a^2
\tilde{F}}{ d\upsilon }
\right]
=
a^3
\frac{d \tilde{F}}{ d\upsilon }
a^3
\tilde{F}''
\end{equation*}

Similarly for the third derivative:

\begin{equation*}
F'''(\gamma)= 
\frac{d }{ d\gamma }
\left[
\frac{d^2 F}{ d^2\gamma }
\right]
=
a^4
\frac{d \tilde{F}}{ d\upsilon }
=
a^4
\tilde{F}'''
\end{equation*}

then, substitute into the momentum equation:

\begin{equation} 
F'''
+
F
F''
=
a^4
\tilde{F}'''
+
a^4
\tilde{F}
\tilde{F}''
=
a^4
[
\tilde{F}'''
+
\tilde{F}
\tilde{F}''
]
=
0
\end{equation} 

Then, $\tilde{F}$ the same equation that $F$.

To apply the invariance of the equation, consider the boundary condition at $\infty$.

\begin{equation*}
F'_\infty
=
a^2
\tilde{F}'_\infty
=1
\end{equation*}

Then, the its not necessary to shooting the value of $\tilde{F}_\infty'$
because any value is solution, only need to be rescaled using the factor a,  
\begin{equation*}
a
=\sqrt{\frac{1}{\tilde{F}'_\infty}}
\end{equation*}

To resolve we are using the reduction of order
defining

\begin{equation}
    S=
\begin{bmatrix}
   S_1\\ 
   S_2\\
   S_3\\
\end{bmatrix}
=
\begin{bmatrix}
   F\\ 
   F'\\
   F''\\
\end{bmatrix}
,
\quad
S'=
\begin{bmatrix}
   F'\\ 
   F''\\
   F'''\\
\end{bmatrix}
=
\begin{bmatrix}
   F'\\ 
   F''\\
   -FF''\\
\end{bmatrix}
=
\begin{bmatrix}
S_2\\
S_3\\
-S_1S_3\\
\end{bmatrix}
\end{equation}

Then the stream-wise equation gets  

\begin{equation}
S'=
\begin{bmatrix}
   F'\\ 
   F''\\
   F'''\\
\end{bmatrix}
=
\begin{bmatrix}
S_2\\
S_3\\
-S_1S_3\\
\end{bmatrix}
=
\begin{bmatrix}
0    & 1 & 0\\
0    & 0 & 1\\
-S_3 & 0 & 0\\
\end{bmatrix}
\begin{bmatrix}
   S_1\\ 
   S_2\\
   S_3\\
\end{bmatrix}
\end{equation}

Once the function $F$ is found, 
the $Z$ can be resolved. 
The special form of the species equation

\begin{equation*}
    Z''+FZ'=0 ,
\end{equation*}

with the  boundary conditions 
$Z(\infty):=Z_\infty=Y_F$ and $Z(-\infty):=Z_{-\infty}=-Y_O/n$. 
is also self-similar respect 
to the invariant transformation $\tilde{Z}=aZ(\gamma)+b$. 

Using the chain rule the 
first derivative  

\begin{equation}
   Z'
   =
   \frac{d Z}{d \gamma} 
   =
   \dfrac{d \left[\dfrac{(\tilde{Z}-b)}{a}\right]}{d \gamma} 
   =
   \frac{1}{a}
    \dfrac{d\tilde{Z}}{d \gamma} 
    =
   \frac{1}{a}
    \tilde{Z}'
\end{equation}

and the second derivative

\begin{equation}
   Z''
   =
   \frac{d }{d \gamma} 
   \left[
    \frac{d Z}{d \gamma} 
   \right]
   =
   \frac{1}{a}
   \frac{d }{d \gamma} 
   \left[
    \dfrac{d\tilde{Z}}{d \gamma} 
   \right]
    =
   \frac{1}{a}
    \tilde{Z}''
    .
\end{equation}

Substituting in the $Z$ equation:
can be written in similar form using the reduction of order  

\begin{equation}
    Z''
    +
    Z' 
    =
   \frac{1}{a}
    \tilde{Z}''
    +
   \frac{1}{a}
    F\tilde{Z}'
    =
    \tilde{Z}''
    +
    F\tilde{Z}'
    =
    0
\end{equation}

Then, the equation is self-similarity and the scaling factors
can be obtained using the boundary conditions. 

Beginning with the lower boundary:

\begin{equation}
    Z_{-\infty}=-\frac{Y_O}{n}=a\tilde{Z}_{-\infty} +b .
\end{equation}

and  assuming that $\tilde{Z}_{-\infty}=0$ as a first guess for the shooting method

\begin{equation}
b=-\frac{Y_O}{n} .
\end{equation}

Then, in the upper boundary

\begin{equation}
Z_\infty
=
Y_F=a\tilde{Z}_\infty +b 
=
Y_F=a\tilde{Z}_\infty
-
\frac{Y_O}{n} .
\end{equation}

then

\begin{equation}
a
=
\dfrac{Y_F+\dfrac{Y_O}{n}}{\tilde{Z}_\infty}
.
\end{equation}

To resolve, the reduction of 
order will  be used, defining 

\begin{equation}
    S=
\begin{bmatrix}
   S_1\\ 
   S_2\\
\end{bmatrix}
=
\begin{bmatrix}
   Z\\ 
   Z'\\
\end{bmatrix}
,
\quad
S'=
\begin{bmatrix}
   Z'\\ 
   Z''\\
\end{bmatrix}
=
\begin{bmatrix}
   Z'\\ 
   -FZ'\\
\end{bmatrix}
=
\begin{bmatrix}
S_2\\
-FS_2\\
\end{bmatrix}
\end{equation}

Then the $Z$ equation gets  

\begin{equation}
S'=
\begin{bmatrix}
   Z'\\ 
   Z''\\
\end{bmatrix}
=
\begin{bmatrix}
S_2\\
-FS_2\\
\end{bmatrix}
=
\begin{bmatrix}
0    & 1 \\
0    & -F \\
\end{bmatrix}
\begin{bmatrix}
   S_1\\ 
   S_2\\
\end{bmatrix}
\end{equation}

As the total energy has the same 
form as the $Z$ equation, the same 
procedure is apply. The 
boundary condition are determined by:

\begin{equation}
    H
    =
    \sum_i^NY_ic_{pi}T
    +
    \sum_i^NY_ih_{fi}^{\circ }
    +
    \frac{U^{2}}{2}  
    =
    \sum_i^NY_ic_{pi}T
    \frac{R_i}{R_i}
    +
    \sum_i^NY_ih_{fi}^{\circ }
    +
    \frac{U^{2}}{2}  
    =
    \sum_i^NY_iR_iT
    \frac{\gamma_i}{\gamma_i-1}
    +
    \sum_i^NY_ih_{fi}^{\circ }
    +
    \frac{U^{2}}{2}  
\end{equation}
evaluated in the boundaries $H(\infty): =H_\infty$ and 
$H(-\infty):=H_{-\infty}$
Remember that:

\begin{equation}
M_w=\sum_i^NY_iM_{wi}, 
\quad
\text{with}
\quad
R_i=\frac{R_u}{M_w},
\quad
\text{then}:
\quad
R=\sum_i^NR_iW_i
\end{equation}

so,

\begin{equation}
    H
    =
    \sum_i^NY_iT
    R
    \frac{\gamma_i}{\gamma_i-1}
    +
    \sum_i^NY_ih_{fi}^{\circ }
    +
    \frac{U^{2}}{2}  
    .
\end{equation}

 Then, in the upper boundary

\begin{equation}
    H_\infty
    =
    R_FT_\infty
    \frac{\gamma_F}{\gamma_F-1}
    +
    Y_Fh_{fF}^{\circ }
    +
    \frac{U_\infty^{2}}{2}  
    , 
\end{equation}

and in the lower boundary

\begin{equation}
    H_{-\infty}
    =
    R_OT_{-\infty}
    \frac{\gamma_O}{\gamma_O-1} +
    Y_Oh_{fO}^{\circ }
    +
    \frac{U_{-\infty}^{2}}{2}  
    .
\end{equation}

With definition of $Z_1$, $Z_1$ and $Z_3$, 

\begin{equation}
   Z_1=Y_F-\frac{Y_O}{n}, 
   \quad
   Z_2=Y_F+\frac{Y_p}{n+1}, 
   \quad
   Z_3=\frac{Y_O}{n}+\frac{Y_p}{n+1}, 
\end{equation}

It is important to note that $Z_3=Z_2-Z_1$, which show that 
the tree equation are not independent.

The species equation 

\begin{equation}
   Y_F+Y_O+Y_P=1 
\end{equation}

can be rewritten in terms of $Z_i$ as 

\begin{equation}
   %Y_F
   Z_1
   +
   \frac{Y_O}{n} 
   +
   Y_O
   +
   (n+1)Z_3
   -
   Y_O\frac{n+1}{n}, 
   =
   1 
\end{equation}

\begin{equation}
   Z_1
   +
   Y_O
   \left(
   \frac{1}{n} 
   +
   1
   -
   \frac{n+1}{n}, 
   \right)
   +
   (n+1)Z_3
   =
   1 
\end{equation}

\begin{equation}
   Z_1
   +
   Y_O
   \left(
   \frac{2}{n} 
   \right)
   +
   (n+1)Z_3
   =
   1 
\end{equation}

Resolving for $Y_O$

\begin{equation}
   Y_O
   =
   \frac{n}{2} 
   \left(
   1 
   -
   (n+1)Z_3
   -
   Z_1
   \right)
\end{equation}

and using to found $Y_F$ 

\begin{equation}
   Y_F
   =
   Z_1
   +
   \frac{Y_O}{n}, 
   =
   Z_1
   +
   \frac{1}{2} 
   \left(
   1 
   -
   (n+1)Z_3
   -
   Z_1
   \right)
   =
   \frac{1}{2} 
   Z_1
   +
   \frac{1}{2} 
   \left(
   1
   -
   (n+1)Z_3
   \right)
\end{equation}

Now, for the $Z_3$

\begin{equation}
   Y_p 
   =
   (n+1)
   \left(
       Z_2
       -
       Y_F
   \right)
   =
   (n+1)
   \left(
       Z_2
       -
   \frac{1}{2} 
   Z_1
   -
   \frac{1}{2} 
   \left(
   1
   -
   (n+1)Z_3
   \right)
   \right)
\end{equation}

So, the mass fraction can be expressed in terms of the $Z_i$ as expected.  However, as $C_p$ was considered constant is more adequate use it to 
calculate the temperature.

Temperature is calculated in function of the $H$:

\begin{equation}
T
=
\frac{\left(    
    H(\gamma)
    -
    Y_F
    \Delta^{\circ}_{fF}
    -
    Y_O
    \Delta^{\circ}_{fO}
    -
    \dfrac{u^2}{2}
\right)}
{C_pZ(\gamma)}
\end{equation}

    {\color{red}For mim is not consistent to have the formation enthalpy of the products equal to zero}

with 

\begin{equation}
C_p(Z(\gamma))
=
\sum_i^N
Y_i
C_{pi}(\gamma)
.
\end{equation}

To calculate the $Y_i$ is used the variable $Z$ once was assumed
an infinite reaction rate, fast chemistry. Then the oxygen and fuel can not presented at the same place ate the same time.  Remember that $Z_1=Y_F-Y_O/n$.
Then, the mass fraction can be calculated using the function $max$ and $min$ defined as  

\begin{equation}
Y_F=
[\max(Z_1(\gamma), 0] 
= 
\begin{cases} Z_1(\gamma) & \text{if } Z_1(\gamma) \geq 0 \\ 0 & \text{if }  Z_1(\gamma) < 0 \end{cases}    
\end{equation}

and 

\begin{equation}
Y_O=
[-n\min(Z_1(\gamma), 0]
= 
\begin{cases} Z_1(\gamma) & \text{if } Z_1(\gamma) \leq 0 \\ 0 & \text{if }  Z_1(\gamma) > 0 
\end{cases}    
\end{equation}

Then, $Yp=1-Y_F-Y_O$, and the temperature can be calculated. 

The heat release can be written as :

\begin{equation*}
    \Delta H_f^{\circ} =\sum_{products}^Ns_i\Delta h_{fi}^{\circ}-\sum_{reactans}^Ms_i\Delta h_{fi}^{\circ}    
\end{equation*}

For a exothermic reaction, the heat release  $Q=-\Delta H_f^{\circ}$  with $\Delta H_f^{\circ}<0$, then the flame temperature can be defined as 

\begin{equation}
    T_{ad}
    =
    \frac{Q}{C_p}+T_u
\end{equation}

when $T_{ad}$ is the adiabatic temperature of the flame and $T_u$ is the unburned temperature of the reactants. 

For one step reaction $F+nO\rightarrow (n+1)P$ the energy balance is 

\begin{equation*}
    \Delta h_{fF}^{\circ}
    +
    n
    \Delta h_{fF}^{\circ}
    +
    C_{p,reactans}
    [T_u-T^{\circ})]
    =(n+1)[h_{fP}^{\circ}+C_{pP}(T_{ad}-T^{\circ})]
\end{equation*}

where $T^{\circ}$ is the reference temperature at which the standard formation
enthalpies are defined. Usually $T^{\circ}=298.15 K$